% DMRN+21: the official 2026 kit ships the 2025 style; its event header and
% footer are updated in the local copy, without changing its page layout.
\documentclass[a4paper]{article}
\usepackage[T1]{fontenc}
\usepackage{caption}
\usepackage{times,amsmath,graphicx,url,dmrn2026}
\usepackage{indentfirst}
\usepackage{xcolor}
\captionsetup[figure]{skip=3pt}\setlength{\intextsep}{4pt plus 1pt minus 1pt}
\setlength{\parskip}{0.1em}
\captionsetup[figure]{font=footnotesize,labelfont=footnotesize}
\captionsetup[table]{font=footnotesize,labelfont=footnotesize}
\setlength{\belowcaptionskip}{0pt}
\renewcommand{\thesection}{\Roman{section}}
\newcommand{\head}[1]{\par\vspace{2pt}\noindent\textbf{#1.}\ \ignorespaces}
\providecommand{\bstctlentry}[1]{\nocite{#1}}
% Layout economies after the DMRN+15 auraloss one-pager: tighter title block,
% 7 pt footer close to the page edge, taller text block. The style file is unchanged.
\setlength{\textheight}{266truemm}\setlength{\footskip}{0.2in}
\cfoot{\fontsize{7}{8}\selectfont\textsc{DMRN+21: DIGITAL MUSIC RESEARCH NETWORK ONE-DAY WORKSHOP 2026, KING'S COLLEGE LONDON, FRI 11 DECEMBER 2026}}
\makeatletter
\def\@maketitle{\newpage
 \null
 \begin{minipage}{0.5\textwidth}
 \textsc{DMRN+21: Digital Music Research Network\\One-day Workshop 2026}\vspace{3mm}\\\textsc{King's College London}\\\textsc{Fri 11 December 2026}
 \end{minipage}%
 \begin{minipage}{0.5\textwidth}
 \hfill\hbox{\includegraphics[width=6cm]{DMRN_Logo_high_res.jpg}}
 \end{minipage}
 \vskip 1em \begin{center}
 {\large \bf \@title \par} \vskip 0.8em {\large
\begin{tabular}[t]{c}\@name \\ \@address
 \end{tabular}
 \par}
 \end{center}
 \par
 \vskip 0.6em}
\makeatother

\title{Efficient Backprop for an FDTD Kirchhoff--Carrier String}
\name{\normalfont Pablo Tablas de Paula, Emmanouil D. Benetos and Joshua D. Reiss\thanks{This work was supported by UKRI grant EP/S022694.}}
\address{\small Centre for Digital Music, Queen Mary University of London, p.tablasdepaula@qmul.ac.uk}

\begin{document}
\bstctlentry{BSTcontrol}
\ninept
\maketitle
\begin{sloppy}
\begin{abstract}
We present a differentiable FDTD Kirchhoff--Carrier string whose
backpropagation runs in a custom CUDA kernel that stores only periodic
checkpoints and three scalars per time step. A fourth-order spatial stencil
removes grid dispersion, allowing lower internal oversampling. Peak memory
stays nearly flat with batch size, and gradients are faster and far leaner
than with a state-of-the-art differentiable modal solver.
\end{abstract}

\section{Introduction}
Although FDTD synthesis is increasingly leaner~\cite{bilbao2024guitar},
differentiable frameworks~\cite{engel2020ddsp} largely keep away from it,
since backpropagating through a time-stepping scheme results in deep
recurrent differentiation graphs, expensive in time and memory. Instead, the
literature has favoured cheaper approaches like digital
waveguides~\cite{tablas2026waveguides,tablas2026sound}, which do not
naturally capture nonlinearities and collisions, and modal
synthesis~\cite{diaz2025}. In this work we demonstrate that an FDTD string
can be backpropagated faster, and in far less memory, than a current
differentiable modal solver.

\section{Model and method}
The string is the damped Kirchhoff--Carrier stiff string, simply supported at
both ends:
\begin{equation}
\begin{aligned}
\rho A\,u_{tt}&=T_0\,G\,u_{xx}-EIu_{xxxx}-2\rho A\sigma_0u_t+2\rho A\sigma_1u_{txx},\\
G&=1+\frac{EA}{2LT_0}\int_0^L\!u_x^2\,dx,
\end{aligned}
\label{eq:pde}
\end{equation}
with density $\rho$, cross-section $A$, Young's modulus $E$, second moment of
area $I$, tension $T_0$ and damping $\sigma_0,\sigma_1$; the stretch of the whole
string raises its tension by the factor $G$, coupling every mode through one
integral.

\head{Time stepping}
We use Russo et al.'s scheme~\cite{russo2026} in Bilbao's
notation~\cite{bilbao2009}: $u^n_l\approx u(lh,nk)$ with grid spacing $h$ and
time step $k$; $\delta_{tt}$, $\delta_{t\cdot}$, $\delta_{t-}$ are the second,
centred and backward time differences, $\mu_{t\cdot}u^n=\tfrac12(u^{n+1}+u^{n-1})$
the centred time average, $\delta_{x+}$, $\delta_{xx}$, $\delta_{xxxx}$ the
spatial differences and $\langle f,g\rangle=h\sum_lf_lg_l$ the inner product on
the grid:
\begin{equation}
\rho A\,\delta_{tt}u=T_0\,g\,\delta_{xx}u-EI\delta_{xxxx}u
 -2\rho A\sigma_0\delta_{t\cdot}u+2\rho A\sigma_1\delta_{t-}\delta_{xx}u,
\label{eq:scheme}
\end{equation}
\begin{equation}
g^n=1+\frac{EA}{2LT_0}\,\big\langle\delta_{x+}u^n,\;\delta_{x+}\mu_{t\cdot}u^n\big\rangle.
\label{eq:g}
\end{equation}
The tension-modulation factor $g^n$ carries the whole nonlinearity: it is $G$
with the stretch averaged over the two neighbouring time levels, the choice that
conserves the discrete energy~\cite{bilbao2009}. Because it reaches $u^{n+1}$
the scheme is implicit, but only through this one scalar, which is linear in
$u^{n+1}$: evaluating (\ref{eq:g}) with the linear prediction $\tilde u^{n+1}$
(the update with $g=1$) in place of $u^{n+1}$ and dividing $g^n-1$ by
$1+\frac{k^2EA}{4L\rho A(1+\sigma_0k)}\|\delta_{xx}u^n\|^2$ gives $g^n$ exactly
(Russo et al.'s Sherman--Morrison step). In the discrete sine basis every spatial
operator is diagonal, so with $g^n$ known (\ref{eq:scheme}) updates each mode
independently.

\head{Cached adjoint}
Instead of every state of every step, we keep $(u^n,u^{n-1})$ every 256 steps
and, per step, only the three scalars behind $g^n$: $g^n$ itself,
$\|\delta_{xx}u^n\|^2$ and $\langle\delta_{x+}u^n,\delta_{x+}\tilde u^{n+1}\rangle$.
On the reverse pass each segment is re-run from its checkpoint, cheaply since
with $g^n$ known (\ref{eq:scheme}) is diagonal, and differentiated backwards,
the stored scalars supplying the derivatives of $g^n$ through (\ref{eq:g}). A
CUDA kernel runs this in parallel over modes and batch items.

\head{Spatial stencil}
Spatial derivatives use the fourth-order five-point
stencil~\cite{bilbao2009}, closed at the simply supported ends by odd
reflection; it is diagonal in the same sine basis, so only the mode
eigenvalues change, and on our 176-interval grid the waveform error against
a converged solution of Russo et al.'s scheme~\cite{russo2026} drops from
11.7\% to 1\% at 2$\times$ oversampling.

\section{Evaluation}
Each evaluation is one sound-matching step: a batch of strings is simulated
for 1 s, the mean absolute error of each waveform against the target's is
computed, and backpropagated to ten parameters: the eight physical ones of
(\ref{eq:pde}) plus the pluck and pickup positions. The target is Russo et
al.'s steel string ($\rho=8050$ kg/m$^3$, $T_0=75$ N, $r=0.356$ mm, $E=174$ GPa,
$L=1$ m, $\sigma_0=0.92$ s$^{-1}$, $\sigma_1=2.9\times10^{-4}$ m$^2$/s), plucked
at its centre with a 2 mm raised cosine of half-width $L/4$ and read at its
midpoint; the batch perturbs every parameter by up to $\pm3\%$. We compare our
method with direct PyTorch autodiff of the same FDTD and a 100-mode
\texttt{jaxdiffmodal} Kirchhoff--Carrier string~\cite{diaz2025} on an A100
80 GB in float64; the FDTD uses 176 intervals at 88.2 kHz, decimated to
44.1 kHz. As seen in
Fig.~\ref{fig:benchmark}, our method is nearly three orders of magnitude
faster than plain autodiff and faster than \texttt{jaxdiffmodal}, and its
memory grows slowly with batch size, whereas both baselines grow linearly and
eventually run out.

\begin{figure}[!h]
 \centering
 \includegraphics[width=\columnwidth]{figures/benchmark-dots-label.pdf}
 \caption{Raised-cosine pluck, 1 s, ten gradients; numbers: time.}
 \label{fig:benchmark}
\end{figure}

\bibliographystyle{IEEEtran}
{\footnotesize\bibliography{references}}
\end{sloppy}
\end{document}
